\chapter{Life Cycle of Cyclones in the Southwestern Atlantic}\label{ch:life_cycle}

This chapter delves into the varied life cycle configurations, statistical characteristics, and geographical distribution of cyclones within the Southwestern Atlantic, with a particular focus on the South American Southeastern region (SESA). Employing the Cyclophaser program, this analysis facilitates a detailed examination of cyclone behaviors and patterns. While some of the results discussed herein are also presented in \citep{deSouza2024}, this chapter narrows its focus to the SESA region, unlike the cited study which includes all cyclogenesis regions across the South Atlantic, extending to the Antarctic Peninsula and the Weddell Sea. Section \ref{sec:climatology_statistics} examines statistical metrics for the selected cyclonic systems' whole life cycle, as well as for individual development phases, while Section \ref{sec:spatial_distribution} explores the spatial distribution of the distinct phases across the SESA and adjacent regions. These two approaches combined not only provide a novel perspective on the South Atlantic storm tracks and on the development of cyclonic systems but also open avenues for future research, as will be discussed.

\section{Climatology and Statistics}
\label{sec:climatology_statistics}

This section delves into the climatology and statistical characteristics of cyclones in the Southwestern Atlantic, focusing on their frequency, seasonality, and developmental metrics such as average duration and intensity. Previous analyses have explored these aspects for systems originating in the SESA region; here, they serve both for comparative purposes and for validating the database generated during this study. A novel aspect of this analysis is the dissection of these metrics not just over the cyclones’ complete lifecycles but also across distinct developmental phases. Section \ref{sec:climatological_aspects} presents the overall characteristics and life cycle configurations of the detected cyclonic systems, Section \ref{sec:cyclone_statistics} explores lifecycle statistics for different genesis regions and seasons, and Section \ref{sec:life_cycle_phase_statistics} provides detailed statistical breakdowns by individual development phases.


\subsection{Climatological Aspects and Life Cycle Configurations}
\label{sec:climatological_aspects}

Using the tracking methodology described in Section \ref{track_method}, the TRACK algorithm identified 33,376 cyclone systems in the South Atlantic region from 1979 to 2020. The focus here was on systems originating from the SESA region, specifically from SE-BR, LA-PLATA, and ARG genesis regions. After filtering to include only these tracks, the number of systems was narrowed down to 7,931 (Figure \ref{fig:pie_climatology}). Among these, 56\% originated in the ARG region, 23.6\% in LA-PLATA, and 20.4\% in SE-BR. This distribution and general seasonality aligns with previous studies by \citet{gramcianinov2019properties} and \citet{crespo2021potential}. However, the SE-BR region exhibited a higher relative frequency in this analysis, potentially attributable to the finer resolution of the ERA5 dataset used here compared to previous studies \citep{gramcianinov2020analysis}. However, comparing the seasonal distribution with earlier studies introduces complexities due to variations in the definitions and boundaries of genesis regions \citep[e.g.,]{reboita2010regimes, crespo2021potential}.

\begin{figure}[ht]
    \centering
    \includegraphics[width=\textwidth]{figs_4/pie_systems_count.png}
    \caption[Cyclogenesis Climatology in the SESA Region]{Number of systems and their relative frequencies by genesis region in the SESA region for the period 1979-2020. The inner pie chart displays the total frequency of systems by region, while the outer ring shows the seasonal distribution of cyclogenesis within each region.}
\label{fig:pie_climatology}
\end{figure}


During the period from 1979 to 2020, cyclonic systems originating from the ARG, LA-PLATA, and SE-BR regions demonstrated 41 distinct life cycle configurations (Figure \ref{fig:all_lifecycle_configurations}). Most configurations did not account for more than 1\% of the total system count, typically lacking a mature stage. This phenomenon is likely due to these systems having maturation stages shorter than the CycloPhaser's threshold for detecting this phase. Adjusting the threshold to recognize shorter mature stages would risk identifying spurious cycles.


\begin{figure}[ht]
    \centering
    \includegraphics[width=\textwidth]{figs_4/combined_barplots_total_all_systems.pdf}
    \caption[Life Cycle Configurations]{Distribution of all identified cyclone lifecycle configurations in the SESA region, spanning from 1979 to 2020. The configurations are sorted by their total count (frequency), with each bar representing the total number of cyclones exhibiting a specific lifecycle pattern, containing distinct stages: incipient (Ic), intensification (It), mature (M), decay (D), and residual (R), showing their sequence within the lifecycle.}
\label{fig:all_lifecycle_configurations}
\end{figure}

To focus on significant life cycle patterns, it was filtered the configurations to include only those comprising at least 1\% of all types (Figure \ref{fig:filtered_lifecycle_configurations}). This criterion reduced the sample size to 7,531 systems, accounting for approximately 95\% of the original database. It was also merged life cycle counts that differed solely by the inclusion of a residual stage, as these are not indicative of the system's actual development (as discussed in Section \ref{sec:cyclophaser_description}). Further, to avoid bias in phase statistics, it was excluded life cycles lacking a mature phase. For instance, systems exhibiting only intensification and/or decay phases might have longer durations for these phases compared to others, potentially skewing statistics. Such configurations might represent undeveloped systems exiting the tracking domain before maturation. With these filters applied, the number of systems analyzed was reduced to 7,151, representing about 90\% of the initial count of systems with genesis in the SESA region (Figure \ref{fig:filtered_lifecycle_configurations}b).

\begin{figure}[ht]
\centering
\includegraphics[width=\textwidth]{figs_4/combined_barplots_filtered.png}
\caption[Filtered Life Cycle Configurations]{(a) Lifecycle configurations of cyclones that represent at least 1\% of the total identified patterns. (b) Adjusted lifecycle configurations after merging periods classified as residual and removing stages not contributing to a standard development pattern, focusing analysis on the core developmental stages of cyclones.}
\label{fig:filtered_lifecycle_configurations}
\end{figure}

Figure \ref{fig:representative_life_cycles} depicts representative cases for each life cycle configuration retained for analysis (as shown in Figure \ref{fig:filtered_lifecycle_configurations}b). The most common configuration in the SESA region typically follows the expected development pattern of extratropical cyclones: incipient, intensification, mature, and decay stages (Figure \ref{fig:representative_life_cycles}a), accounting for 68\% of all evaluated systems. The second most common type mirrors this sequence but omits the incipient stage, indicating a rapid intensification process (Figure \ref{fig:representative_life_cycles}b). The third and fourth most frequent types depict a repeated complete cycle, including a secondary intensification, mature, and decay phases, with the third type including an incipient stage (Figure \ref{fig:representative_life_cycles}c) and the fourth lacking one (Figure \ref{fig:representative_life_cycles}d). The fifth and sixth configurations exhibit an early decay stage, differentiated by the presence (Figure \ref{fig:representative_life_cycles}e) or absence (Figure \ref{fig:representative_life_cycles}f) of an incipient phase.

\begin{figure}[h!]
\centering
\includegraphics[width=\textwidth]{figs_4/main_species_life-cycle.png}
\caption[Cyclone Life Cycles - Representative Examples]{Representative examples of cyclone life cycles for configurations shown in Figure \ref{fig:filtered_lifecycle_configurations}b. Panels A to F illustrate the central relative vorticity \(\zeta_{850}\) and its first (\(\zeta_{fs}\)) and second (\(\zeta_{fs^2}\)) smoothed derivatives. Background colors indicate different life cycle phases: : incipient (Ic), intensification (It), mature (M) and decay (D). Lines represent the original vorticity series (\(\zeta\)), the first (\(\zeta_{fs}\)), and the second (\(\zeta_{fs^2}\)) smoothed relative vorticity series.}
\label{fig:representative_life_cycles}
\end{figure}

The analysis reveals minimal to negligible seasonal variability in the frequency of each life cycle configuration across different genesis regions (Figure \ref{fig:lifecycle_configurations_seasonal_variability}). The configurations "Ic, It, M, D" and "Ic, It, M, D, It2, M2, D2" are consistently more frequent, jointly accounting for 60\% to 68\% of all cases across all regions. However, exceptions include the "It, M, D" configuration emerging as the second most frequent in ARG during JJA, and in SE-BR during MAM and JJA. Additionally, the "Ic, D, M, D2" configuration appears significantly more often in LA-PLATA during JJA and in SE-BR during SON.


\begin{figure}[h!]
\centering
\includegraphics[width=\textwidth]{figs_4/multi_panel_barplots_filtered.png}
\caption[Life cycle configurations by regions and season]{Similar to Figure \ref{fig:filtered_lifecycle_configurations}, but separated by genesis region (ARG, LA-PLATA and SE-BR) and season.}
\label{fig:lifecycle_configurations_seasonal_variability}
\end{figure}


\subsection{Mean Cyclone Statistics}
\label{sec:cyclone_statistics}

This section explores the statistical analysis of cyclones in the SESA region, comparing the constrained data utilized in this study with extant literature. Results are presented both for an aggregate view across all genesis regions and individually by region. To capture the seasonal variability of cyclone behavior, analysis is focused on the winter (JJA) and summer (DJF) months. This approach is chosen because the transitional seasons (MAM and SON) generally display characteristics that are intermediate to those observed in JJA and DJF, aligning with findings from previous climatologies of the SESA region \citep[e.g.]{gan1991surface, reboita2010south, crespo2021potential}. 

On average, cyclones exhibit longer lifespans during DJF, with a duration of $97.6 \pm 66.2$ hours (Figure \ref{fig:pdf_total_time_total}a) compared to $91.1 \pm 59.7$ hours during JJA (Figure \ref{fig:pdf_total_time_total}b). Conversely, the average total distance traveled by cyclones is slightly greater during JJA, amounting to $4777 \pm 3182$ km (Figure \ref{fig:pdf_total_time_total}d), versus $4658 \pm 3212$ km in DJF (Figure \ref{fig:pdf_total_time_total}c). This is reflected in the mean propagation speed, which is higher in JJA at $15.3 \pm 4.8$ m/s (Figure \ref{fig:pdf_total_time_total}f) compared to DJF at $14.0 \pm 4.4$ m/s (Figure \ref{fig:pdf_total_time_total}e).

\begin{figure}[h!]
\centering
\includegraphics[width=\textwidth]{figs_4/pdf_total_phase_all_metrics.png}
\caption[PDF - Metrics for Each Region and Season]{Probability density functions (PDF) comparing cyclone lifecycle metrics across different genesis regions (ARG, LA-PLATA, SE-BR) for DJF and JJA. Each panel displays the density distributions for a specific metric: total time (A and B), total distance (C and D), mean speed (E and F), mean vorticity (G and H), and mean growth rate (I and J). The lines in each graph represent the mean values for all regions combined (black), ARG (blue), LA-PLATA (green), and SE-BR (red), while the annotations represent the mean and standard deviation values.}
\label{fig:pdf_total_time_total}
\end{figure}

These results largely corroborate those by \citet{gramcianinov2019properties, gramcianinov2020analysis, hoskins2005new} and exceed the mean duration, traveled distance, and propagation speed reported by \citet{simmonds2000mean}, \citet{mendes2010climatology}, and \citet{reboita2010south}. The disparities likely stem from differing tracking methodologies, data sets, and the exclusion of systems without a mature phase in our analysis. Particularly, the exclusion of continental systems may inflate displacement statistics during austral summer when quasi-stationary lows are prevalent over the continent \citep{mendes2010climatology}. Furthermore, \citet{reboita2010south} focused on oceanic systems, which might have overlooked early development stages, and their use of relative vorticity at 10m, impacted by surface drag, could account for variations in propagation speed.

The analysis reveals notable inter-regional variability in cyclone behavior, with trends predominantly mirroring those observed in the ARG region, which accounts for the highest frequency of cyclone occurrences within the SESA region (Figure \ref{fig:pie_climatology}). When examining total cyclone duration, distinct seasonal patterns emerge: both the ARG and SE-BR regions exhibit significantly shorter lifespans during JJA compared to DJF, with the probability density functions (PDFs) peaking at approximately 50 hours. In contrast, the LA-PLATA region shows slightly longer durations during JJA, with a mean difference of about 1.5 hours. Additionally, the distribution of cyclone durations in the LA-PLATA region is more right-skewed than in ARG and SE-BR, indicating greater variability in the lifespan of cyclones originating from LA-PLATA. Across all regions, the standard deviation decreases from DJF to JJA, suggesting a higher occurrence of longer-lasting cyclones during the summer months.

In terms of displacement, the ARG region displays shorter average distances during DJF compared to JJA, whereas LA-PLATA and SE-BR exhibit increases, particularly notable in SE-BR. Similar to the trends in duration, the displacement patterns in ARG and SE-BR exhibit distinct peaks at approximately 2500 km, while LA-PLATA shows a more right-skewed distribution, especially during JJA, indicating greater variability in cyclone travel distances. These regional differences in displacement are also mirrored in the average propagation speeds of the cyclones: during DJF, cyclones in ARG generally move faster than those in LA-PLATA and SE-BR. This changes in JJA, with SE-BR exhibiting the highest increases in speed, making it the region with the most mobile cyclones during this season, followed by ARG and then LA-PLATA. This variation is reflected in the PDFs for each region: during DJF, SE-BR's propagation speed peaks at approximately \(10 \, \text{m s}^{-1}\) and ARG at \(13 \, \text{m s}^{-1}\), while LA-PLATA displays a bimodal distribution with peaks at approximately \(10 \, \text{m s}^{-1}\) and \(15 \, \text{m s}^{-1}\). During JJA, the peak for ARG shifts to approximately \(16 \, \text{m s}^{-1}\), and LA-PLATA's distribution changes to peak at around \(14 \, \text{m s}^{-1}\), while SE-BR assumes a bimodal shape, peaking at roughly \(10 \, \text{m s}^{-1}\) and \(16 \, \text{m s}^{-1}\).

The PDFs for the mean vorticity values of the examined systems display a bimodal, right-skewed distribution with peaks near \(2 \times 10^{-5} \, \text{s}^{-1}\) and \(4 \times 10^{-5} \, \text{s}^{-1}\), with a higher occurrence percentage near the first peak and an average value of \(4 \pm 1.68 \times 10^{-5} \, \text{s}^{-1}\) during DJF (Figure \ref{fig:pdf_total_time_total}g). In JJA, the peaks are positioned near \(2 \times 10^{-5} \, \text{s}^{-1}\) and \(5 \times 10^{-5} \, \text{s}^{-1}\), but with a higher occurrence percentage at the second peak, reflecting an overall increase in mean vorticity, especially in the LA-PLATA region (Figure \ref{fig:pdf_total_time_total}h). The ARG region presents a PDF mirroring the mean cyclone behavior, and SE-BR peaks near \(2 \times 10^{-5} \, \text{s}^{-1}\) in both seasons, being more right-skewed in JJA, which reflects the occurrence of weaker systems in this region. Meanwhile, LA-PLATA shifts from being right-skewed in DJF to left-skewed in JJA, reflecting a significant increase in the intensity of systems in this region. These results highlight the previous findings that LA-PLATA presents the most intense systems in the SESA regions, while SE-BR, the weakest \citep{simmonds2000mean, reboita2010south, gramcianinov2019properties}. Although the regional and seasonal variability largely corroborates findings from \citet{gramcianinov2019properties}, the intensity of cyclones reported here is somewhat smaller, possibly due to the ERA5 dataset be able to detect less intense systems \citep{gramcianinov2020analysis}.

The mean growth rate distribution, is symmetric around zero with a low standard deviation, indicating a dynamic equilibrium in system intensity throughout the life cycle, marked by neither consistent intensification nor weakening, which will be later discussed for each individual phase. It is important to note that these values were computed from normalized relative vorticity data, where positive values suggest system intensification (\(\zeta_{850}\) increasing in magnitude), while negative values indicate a decay (\(\zeta_{850}\) decreasing in magnitude). Overall, although the differences are minimal, the results suggest slightly more intense intensification during DJF (Figure \ref{fig:pdf_total_time_total}i) compared to JJA (Figure \ref{fig:pdf_total_time_total}j), with the latter exhibiting higher mean values.

\subsection{Life Cycle Phase Statistics}
\label{sec:life_cycle_phase_statistics}

This section examines the statistical characteristics of cyclones, exploring various metrics describing cyclonic behavior throughout distinct lifecycle phases. This unique approach is enabled by the Cyclophaser program, offering a more granular analysis than previously available in the literature. As shown in Section \ref{sec:cyclone_statistics}, the region mean results predominantly reflect the ARG region's behavior and hence, region-specific results are detailed here.

Figure \ref{fig:pdf_total_time} displays the PDFs for the mean total time spent in each development phase. These PDFs are right-skewed, indicating high variability across different phases but smaller seasonal and regional differences. Typically, the incipient, mature, and second mature phases show the shortest average durations, often peaking between 3 and 5 hours. The intensification, second intensification, and decay phases often exhibit longer tails, suggesting a wide variability in their durations. In the ARG region, all phases show a slight reduction in mean duration from DJF (Figure \ref{fig:pdf_total_time}a) to JJA (Figure \ref{fig:pdf_total_time}b), with decay and second decay stages exhibiting the most significant reductions, indicating that shorter average total durations in JJA are primarily due to these phases (Figure \ref{fig:pdf_total_time_total}). For LA-PLATA and SE-BR, notable changes include particularly variations in the standard deviations of the second intensification and decay phases, and in the decay phase for SE-BR, suggesting that seasonal differences in total mean durations for these regions are largely attributable to fluctuations in extreme values.

\begin{figure}[h!]
\centering
\includegraphics[width=\textwidth]{figs_4/pdf_total_time.png}
\caption[PDF - Total Time]{Probability density functions (PDF) comparing the mean total duration of each cyclone phase, across all regions combined and for each genesis region (ARG, LA-PLATA, SE-BR) for DJF and JJA. Annotations display mean and standard deviation values for each phase.}
\label{fig:pdf_total_time}
\end{figure}

The total traveled distance for each phase is depicted in Figure \ref{fig:pdf_total_distance}, which mirrors the overall behavior observed for total time (Figure \ref{fig:pdf_total_time}). All PDFs are right-skewed, and across all regions and seasons, the incipient, mature, and second mature phases peak approximately between 300 and 400 km, while the other phases present long tail distributions. In the ARG region, there is an overall reduction in traveled distance across all phases from DJF to JJA. Conversely, for LA-PLATA and SE-BR, there is a general increase, reflecting the seasonal behavior for the entire lifecycle's total distance (Figure \ref{fig:pdf_total_time_total}).

\begin{figure}[h!]
\centering
\includegraphics[width=\textwidth]{figs_4/pdf_total_distance.png}
\caption[PDF - Total Distance]{Similar to Figure \ref{fig:pdf_total_time}, but for total traveled distance.}
\label{fig:pdf_total_distance}
\end{figure}

The results for mean speed (Figure \ref{fig:pdf_mean_speed}) aggregate the results for total time and traveled distance. There is less variability across distinct phases for this metric when compared to the total time and distance, with the PDFs for distinct phases, regions, and seasons often presenting a Gaussian-like shape centered between 10 and 15 \(m s^{-1}\). These values align with those found in previous studies \citep{gramcianinov2019properties, gramcianinov2020analysis, hoskins2005new}, although slightly higher, likely due to the constraints adopted in this analysis. Additionally, this study's ability to account for distinct developmental phases provides a foundation for further research into the dynamical mechanisms that could be linked to the systems' varying speeds across different periods. In the ARG region, the largest increases in mean propagation speed from DJF to JJA are found in the secondary development phases, accompanied by increases in standard deviation values, with slightly right-skewed PDFs for these phases. LA-PLATA displays a similar behavior, especially for intensification and mature decay, which shifts their peaks from near \(10 m s^{-1}\) in DJF to near \(15 m s^{-1}\) in JJA. For SE-BR, the major change is in the shape of the PDFs, which change from a Gaussian-like shape in DJF to a more right-skewed shape in JJA. Overall, although the differences are only marginal, the cyclones tend to be faster in the incipient stage and in the intensification stage, being slower at the mature phase, speeding up again in the decay phase. Also, the systems tend to be slower in the second development cycle than in the first one, especially slower at the second mature phase. Therefore, cyclones in the SESA region are slower at their mature phases. The slower propagation speed of cyclones is related to extreme wave occurrence \citep{gramcianinov2023impact} and therefore the positioning of these systems during that phase are of importance for coastal management in the South American Southeastern coast.

\begin{figure}[h!]
\centering
\includegraphics[width=\textwidth]{figs_4/pdf_mean_speed.png}
\caption[PDF - Mean Speed]{Similar to Figure \ref{fig:pdf_total_time}, but for mean propagation speed.}
\label{fig:pdf_mean_speed}
\end{figure}

The PDFs for mean central \(\zeta_{850}\) across various cyclone phases exhibit bimodal and right-skewed distributions, with a general increase in mean vorticity from DJF to JJA (Figure \ref{fig:pdf_mean_vorticity}). This trend reflects the heightened intensity of cyclonic systems during JJA compared to DJF (Figure \ref{fig:pdf_total_time_total}). It is important to note that these vorticity values are derived from the raw output of the TRACK program, and as such, they may include some level of noise that can influence the statistical results. Notably, the mature phase is characteristically the most intense during the first developmental cycle across all regions, demonstrating a wide range of vorticity values and highlighting the variability in peak intensities of mature cyclones. Meanwhile, the second mature phase typically records the highest vorticities among all phases, suggesting a tendency for systems to intensify further during subsequent cycles. In contrast, the incipient phase generally exhibits the lowest mean vorticities, reflecting the nascent stages of cyclone development. However, in the SE-BR region, this phase shows unusually high mean vorticities during JJA, occasionally surpassing other phases, indicating exceptionally conducive conditions for early cyclone formation. The intensification and decay phases generally present similar mean vorticities, aligning with their respective roles in modifying system intensity: intensification increases vorticity from a baseline state, while decay reduces it back towards this baseline. Secondary intensification and decay phases often exhibit higher mean vorticities than their initial counterparts, reflecting the dynamics of re-intensification in the lifecycle, particularly marked by the higher mean values observed during the second mature phase.


\begin{figure}[h!]
\centering
\includegraphics[width=\textwidth]{figs_4/pdf_mean_vorticity.png}
\caption[PDF - Mean Vorticity]{Similar to Figure \ref{fig:pdf_total_time}, but for mean central relative vorticity at cyclone center ($zeta_{850}$).}
\label{fig:pdf_mean_vorticity}
\end{figure}

The results found here for the incipient stage exhibit higher vorticities compared to the initial vorticities reported by \citet{reboita2010south}, \citet{gramcianinov2019properties}, and \citet{gramcianinov2020analysis}. This discrepancy may be attributed to our exclusion criteria, which likely excluded weaker and underdeveloped systems, coupled with the observed tendency for vorticity to increase (becoming more negative) when an incipient stage is followed by an intensification phase (e.g., Figure \ref{fig:representative_life_cycles}a and c), potentially resulting in a bias toward higher values for this phase. Additionally, \citet{sinclair1995climatology}, in their analysis focusing on the first time step, the time step of minimum vorticity, and the last time step, reported higher vorticity values compared to our findings for the incipient, mature, and decay stages. However, these authors omitted the influence of topography by using geostrophic relative vorticity, which justifies the difference in vorticity distribution observed in their study.

There are complexities in interpreting the results for the mean growth rate due to the use of raw vorticity data from the TRACK algorithm, which introduces additional noise to the findings (Figure \ref{fig:pdf_mean_growth_rate}). This metric exhibits considerable variability across different regions, with most PDFs centered around zero yet displaying wide variations, where standard deviations significantly exceed the mean values. The incipient stage often shows the highest or second-highest growth rates, typically skewing more to the right. These observations are consistent with previously reported mean growth rates for cyclogenesis regions by \cite{hoskins2005new} and \cite{gramcianinov2019properties}. A similar pattern is observed in the intensification phase, which alternates with the incipient stage in displaying the greatest growth rates. Despite positive average values, the decay and second decay phases predominantly peak at negative values, reflecting their role in diminishing system vorticity.

\begin{figure}[!htbp]
\centering
\includegraphics[width=\textwidth]{figs_4/pdf_mean_growth_rate.png}
\caption[PDF - Mean Growth Rate]{Similar to Figure \ref{fig:pdf_total_time}, but for mean growth rate.}
\label{fig:pdf_mean_growth_rate}
\end{figure}


While most regions and seasons show the mature phase peaking at zero — expected due to its position near local vorticity minima — the positive mean suggests that the increase in vorticity preceding the peak is more pronounced than the subsequent decline. This assertion is supported by the lower mean values observed during the decay phase compared to the intensification phase. Moreover, the detection of the mature phase often occurs between the peaks and valleys of vorticity derivatives (Section \ref{sec:cyclophaser_description}), potentially shifting its center of mass toward either the intensification or decay stages, as depicted in Figure \ref{fig:mature_phase}. During JJA in LA-PLATA, the mature phase peaks positively, whereas in SE-BR during DJF, it peaks negatively. Determining whether these shifts are due to greater increases in intensity, displacement of the phase's center of mass, or residual data noise is beyond this study's scope.

\begin{figure}[!htbp]
\centering
\includegraphics[width=20pc]{figs_4/mature_phase.pdf}
\caption[Mature Phase - Center of Mass]{Illustrative example of how the center of mass of the mature phase can be displaced: (A) Symmetrically around the vorticity valley (\(\partial \zeta_{valley}\)) when distances to the first derivative valley (\(\partial \zeta_{valley}\)) and peak (\(\partial \zeta_{peak}\)) are equal. (B) Displacement towards the previous phase when \(\partial \zeta_{peak}\) is closer to \(\zeta_{valley}\). (C) Displacement towards the subsequent phase when \(\partial \zeta_{valley}\) is closer to \(\zeta_{valley}\).}
\label{fig:mature_phase}
\end{figure}

\section{Spatial Distribution of Cyclone Life Cycle Phases}
\label{sec:spatial_distribution}

The current section explores the spatial distribution across the SESA region of the distinct phases of cyclone development: incipient (Section \ref{sec:incipient_phase}), intensification (Section \ref{sec:intensification_phase}), mature (Section \ref{sec:mature_phase}), decay (Section \ref{sec:decay_phase}) and the secondary development cycle, as well the residual phase (Section \ref{sec:secondary_development}). The analysis focuses only on systems that presented a mature phase to avoid including underdeveloped systems, those that might not have exhibited a mature phase before leaving the tracking domain, or cases where the maturation stages were shorter than the CycloPhaser's threshold for detecting this phase (as discussed in Section \ref{sec:climatological_aspects}). The track density ranges differ from one figure to the other as does not make sense to compare track density between distinct phases, as those differ in duration and in total displacement (Section \ref{sec:life_cycle_phase_statistics}.


\subsection{Incipient Phase} \label{sec:incipient_phase}

Figure \ref{fig:density_map_incipient} illustrates cyclone track density (cyclones per \(10^6 \, km^2\) per month) during the incipient phase for DJF and JJA in the SESA region. Given that transitional seasons display intermediate characteristics between these extremes \citep[e.g.,]{gan1991surface,reboita2010south, crespo2021potential}, our analysis primarily focuses on DJF and JJA, encompassing the cyclogenesis regions of ARG, LA-PLATA, and SE-BR. As expected, the track density for the incipient phase correlates with these genesis regions (Figure \ref{fig:density_map_incipient}a and b).

\begin{figure}[!htbp]
\centering
\includegraphics[width=\textwidth]{figs_4/density_map_incipient.png}
\caption[Track Density - Incipient Phase]{Cyclone track density during the incipient phase for all systems with genesis in the ARG, LA-PLATA and SE-BR regions. Panels A and B show the aggregate density across all regions for DJF and JJA, respectively. Panels C to H depict individual genesis regions: ARG (C and D), LA-PLATA (E and F), and SE-BR (G and H) for DJF (left panels) and JJA (right panels). The unit is cyclones per \(10^6 \, km^2\) per month.}
\label{fig:density_map_incipient}
\end{figure}

The ARG region consistently shows similar track densities in both DJF (Figure \ref{fig:density_map_incipient}c) and JJA (Figure \ref{fig:density_map_incipient}d), exceeding \(30\) cyclones per \(10^6 \, km^2\) per month. In contrast, LA-PLATA and SE-BR regions exhibit noticeable seasonal patterns. LA-PLATA sees increased activity in JJA (Figure \ref{fig:density_map_incipient}f) compared to DJF (Figure \ref{fig:density_map_incipient}e), with densities rising above \(20\) and \(10\) cyclones per \(10^6 \, km^2\) per month, respectively. Conversely, SE-BR demonstrates higher activity in DJF (Figure \ref{fig:density_map_incipient}g) than in JJA (Figure \ref{fig:density_map_incipient}h), with densities exceeding \(10\) and \(8\) cyclones per \(10^6 \, km^2\) per month, respectively. This seasonality is consistent with findings from previous studies \citep{gan1991surface, hoskins2005new, mendes2010climatology, reboita2010south, gramcianinov2019properties, crespo2021potential}. Variations in reported track density values by \citep{gramcianinov2020analysis, gramcianinov2019properties} may arise from different methodologies; those studies analyzed only the initial time step of the system life cycle, while this study considers all periods classified as the incipient stage. Additionally, \citet{gramcianinov2019properties} utilized the NCEP-CFSR reanalysis database for tracking, in contrast to the ERA5 dataset employed here.

% Summary: overall pattern match the one for gensis regions described in the literature

\subsection{Intensification Phase}  \label{sec:intensification_phase}

The track densities during the intensification phase (Figure \ref{fig:density_map_intensification}) exhibit significantly higher values compared to the incipient stage. This is expected due to the longer duration of the intensification phase, as discussed in Section \ref{sec:cyclophaser_description}. Across all genesis regions, during both DJF (Figure \ref{fig:density_map_intensification}a) and JJA (Figure \ref{fig:density_map_intensification}b), the maximum track density extends from approximately 60$^\circ$W to 20$^\circ$W. The central concentration around 50$^\circ$S, particularly during both seasons, is primarily composed of systems originating in the ARG region (Figure \ref{fig:density_map_intensification}c and d). During DJF, the maximum density region, extending eastward/southeastward from the incipient stage, reaches up to \(100\) cyclones per \(10^6 \, km^2\) per month, while in JJA, it shows a slight decrease to \(80\) cyclones per \(10^6 \, km^2\) per month, with a predominantly eastward orientation. The higher track density and the cyclones reaching further east in DJF can be attributed to the longer mean duration of the intensification phase for systems from ARG, as shown in Figure \ref{fig:pdf_total_time}, and the overall greater displacement distances during this phase across all genesis regions, as indicated in Figure \ref{fig:pdf_total_distance}.

\begin{figure}[!htbp]
\centering
\includegraphics[width=\textwidth]{figs_4/density_map_intensification.png}
\caption[Track Density - Intensification Phase]{Cyclone track density during the intensification phase for all systems with genesis in the ARG, LA-PLATA and SE-BR regions. Panels A and B show the aggregate density across all regions for DJF and JJA, respectively. Panels C to H depict individual genesis regions: ARG (C and D), LA-PLATA (E and F), and SE-BR (G and H) for DJF (left panels) and JJA (right panels). The unit is cyclones per \(10^6 \, km^2\) per month.}
\label{fig:density_map_intensification}
\end{figure}

Moreover, the maximum track density at around \(50^\circ\)S for the ARG region connects with the track densities for the LA-PLATA and SE-BR regions. LA-PLATA shows higher activity during JJA (Figure \ref{fig:density_map_intensification}f), with track densities reaching \(50\) cyclones per \(10^6 \, km^2\) per month, compared to \(40\) during DJF (Figure \ref{fig:density_map_intensification}e). Conversely, SE-BR exhibits higher activity during DJF (Figure \ref{fig:density_map_intensification}g), with densities reaching \(40\) cyclones per \(10^6 \, km^2\) per month, compared to \(20\) during JJA (Figure \ref{fig:density_map_intensification}h). This seasonality for both LA-PLATA and SE-BR regions reflects the seasonality in cyclogenesis (Figure \ref{fig:pie_climatology}) and consequently, of the incipient phase (Figure \ref{fig:density_map_incipient}).  Consequently, in the aggregate densities for all regions during DJF (Figure \ref{fig:density_map_intensification}a), the SE-BR region's signal is more apparent near \(30^\circ\)S and \(40^\circ\)W. In JJA (Figure \ref{fig:density_map_intensification}b), the track density associated with LA-PLATA is evident as the maximum track region extends southeastward from near \(35^\circ\)S and \(55^\circ\)W (Figure \ref{fig:density_map_intensification}f). Furthermore, for the LA-PLATA region, it can be seen that during JJA, the systems tend to reach further southeast, associated with higher mean displacements (Figure \ref{fig:pdf_total_distance}).

Notably, the track densities for the intensification phase across all genesis regions show far-reaching systems that intensify throughout the entire South Atlantic basin, extending as far as \(30^\circ\)E or even \(60^\circ\)E. While the orientation of the track density for the ARG region is predominantly eastward, for LA-PLATA and SE-BR, the tracks are southeastward oriented. This results in a concentration of track densities between \(40^\circ\)S and \(60^\circ\)S, displaced more northward during JJA than DJF, indicating a preference for systems to intensify in this region, which corresponds to the storm track path for the South Atlantic \citep{hoskins2005new, gramcianinov2019properties}.

% Summary: the seasonal differences in the track densities are mostly related to the seasonal differentes in the genesis, with two maxima present in both seasons: one mostly associated with ARG and the other mostly with either LA-PLATA or SE-BR. However, other patterns can be seen: 1) ARG is more eastward-orientated during JJA than DJF and 2) reach further east during DJF; 3)for LA-PLATA the systems tend to reach further southeast during JJA. Overall the tracks a more displaced northward during JJA than DJF.

\subsection{Mature Phase}\label{sec:mature_phase}

In the mature phase, there is a notable reduction in track densities compared to the intensification phase, with maximum values decreasing from \(100\) to to \(20\) cyclones per \(10^6 \, km^2\) per month. During DJF, the aggregate for all genesis regions shows two principal maximum track density areas (Figure \ref{fig:density_map_mature}a). The first area, centered near \(35^\circ\)S and \(45^\circ\)W, exhibits a lower maximum track density reaching \(14\) cyclones per \(10^6 \, km^2\) per month, predominantly associated with cyclones originating in the LA-PLATA and SE-BR regions (Figures \ref{fig:density_map_mature}e and \ref{fig:density_map_mature}g). The second area, located near \(55^\circ\)S and \(30^\circ\)W, reaches up to \(18\) cyclones per \(10^6 \, km^2\) per month and is primarily associated with systems from the ARG region (Figure \ref{fig:density_map_mature}c). During this season, maximum track densities for ARG, LA-PLATA, and SE-BR peak at \(16\), \(9\), and \(10\) cyclones per \(10^6 \, km^2\) per month, respectively.

\begin{figure}[!htbp]
\centering
\includegraphics[width=\textwidth]{figs_4/density_map_mature.png}
\caption[Track Density - Mature Phase]{Cyclone track density during the mature phase for all systems with genesis in the ARG, LA-PLATA and SE-BR regions. Panels A and B show the aggregate density across all regions for DJF and JJA, respectively. Panels C to H depict individual genesis regions: ARG (C and D), LA-PLATA (E and F), and SE-BR (G and H) for DJF (left panels) and JJA (right panels). The unit is cyclones per \(10^6 \, km^2\) per month.}
\label{fig:density_map_mature}
\end{figure}

In JJA, the primary track density maximum for the aggregate is located near \(45^\circ\)S and \(45^\circ\)W, slightly westward compared to the DJF maximum, with secondary maxima to the north and northeast, each reaching \(12\) cyclones per \(10^6 \, km^2\) per month (Figure \ref{fig:density_map_mature}b). This lower density compared to DJF reflects shorter mature phase durations and reduced mobility of systems from the ARG region during JJA, as shown in Figures \ref{fig:pdf_total_time} and \ref{fig:pdf_total_distance}. Notably, the track density center for the ARG region shifts from near \(45^\circ\)W during DJF to around \(55^\circ\)W during JJA(Figure \ref{fig:density_map_mature}d), indicating a tendency for systems to mature closer to their intensification locations, which is also indicated by the lower mobility of these systems when compared to the other genesis regions for this period (Figure \ref{fig:pdf_total_distance}). Secondary maxima near \(35^\circ\)S and \(45^\circ\)S result from overlaps in tracks from all regions, with LA-PLATA contributing significantly (Figure \ref{fig:density_map_mature}f). Although this overlaps with a center from SE-BR, the track density maximum for SE-BR during JJA remains low, not exceeding \(5\) cyclones per \(10^6 \, km^2\) per month(Figure \ref{fig:density_map_mature}h), while for LA-PLATA, it reaches up to \(7\) near \(55^\circ\)W and \(6\) near \(45^\circ\)W.

Overall, the density maximum centers for the mature phase are displaced southeastward from those for the intensification phase for DJF and east/southeastward for JJA. Tracks are typically found between \(40^\circ\)S and \(60^\circ\)S, aligning with the storm track path for the South Atlantic. For DJF, most tracks are located below \(50^\circ\)S, while for JJA, they are predominantly above \(50^\circ\)S. This pattern is more pronounced during the mature phase than during the intensification phase, with the southward shift during DJF largely associated with cyclones originating in the ARG region, while a maximum density region closer to Uruguay and Southern/Southeastern Brazil is mostly linked with LA-PLATA and SE-BR genesis regions.


% Summary: Overall, the maximum track densities are southeast of the intensification centers during DJF and east/southeast during JJA. During JJA, cyclones from ARG region mature close to the intensification. There is a higher positoning cyclones close to Uruguay and Souhtern/Southeastern Brazil during DJF than JJA.

\subsection{Decay Phase} \label{sec:decay_phase}


The decay phase generally exhibits increased track densities compared to the mature phase, with maximum values reaching \(40\) cyclones per \(10^6 \, km^2\) per month, but still lower than observed during the intensification phase. During DJF, two primary centers of maximum track density are observed: one near \(35^\circ\)S, close to Uruguay and South/Southeastern Brazil (Figure \ref{fig:density_map_decay}a), and another near \(55^\circ\)S, west of the Weddell Sea (Figure \ref{fig:density_map_decay}a). The latter, primarily influenced by cyclones from the ARG region, reaches up to \(35\) cyclones per \(10^6 \, km^2\) per month (Figure \ref{fig:density_map_decay}c). The former, influenced by the LA-PLATA and SE-BR regions, records densities of up to \(13\) and \(30\) cyclones per \(10^6 \, km^2\) per month, respectively (Figures \ref{fig:density_map_decay}e and \ref{fig:density_map_decay}g).


\begin{figure}[!htbp]
\centering
\includegraphics[width=\textwidth]{figs_4/density_map_decay.png}
\caption[Track Density - Decay Phase]{Cyclone track density during the decay phase for all systems with genesis in the ARG, LA-PLATA and SE-BR regions. Panels A and B show the aggregate density across all regions for DJF and JJA, respectively. Panels C to H depict individual genesis regions: ARG (C and D), LA-PLATA (E and F), and SE-BR (G and H) for DJF (left panels) and JJA (right panels). The unit is cyclones per \(10^6 \, km^2\) per month.}
\label{fig:density_map_decay}
\end{figure}

For JJA, two primary density maxima are identified: one centered at \(45^\circ\)S, spanning from \(60^\circ\)W to \(20^\circ\)W, and another near \(35^\circ\)S and \(40^\circ\)W, with both areas achieving track densities up to \(35\) cyclones per \(10^6 \, km^2\) per month (Figure \ref{fig:density_map_decay}b). The former is predominantly linked to systems from the ARG region, where densities reach \(30\) cyclones per \(10^6 \, km^2\) per month (Figure \ref{fig:density_map_decay}d), while the latter is influenced by the SE-BR and especially the LA-PLATA regions, showing densities of \(10\) and \(15\) cyclones per \(10^6 \, km^2\) per month, respectively (Figures \ref{fig:density_map_decay}h and \ref{fig:density_map_decay}f). Additionally, a notable center of density, although less pronounced, emerges in the LA-PLATA region, associated with systems that enter an early decay phase (Section \ref{sec:secondary_development}). 

For the ARG region, the decay phase densities are typically shifted southeastward from those of the mature phase during DJF and mostly southeastward/eastward during JJA, with some systems extending westward to the Weddell Sea. This pattern highlights a seasonal variation in the decay behaviors of systems within the SESA region. Moreover, for both LA-PLATA and SE-BR, decay densities are generally proximate to those of the mature phase but extend further southeastward, reaching as far south as \(70^\circ\)S and as far east as \(30^\circ\)E. This expansive spread into the southern Atlantic aligns with findings by \citet{hoskins2005new}, who identify the region south of \(60^\circ\)S and west of the Weddell Sea as the primary dissolution area for cyclones in the South Atlantic.

% Summary: the track density regions tend to elongate southeastward in comparison to the mature phase, with a maximum density east of Weddell sea, related mainly to systems from ARG. Systems with early decay phase appear as density centers in LA-PLATA. During DJF there are more track positioned near Uruguay and South/Southeast Brazil than JJA.

\subsection{Secondary Development Phases}\label{sec:secondary_development}

For analyzing the track densities for the phases related to the systems' secondary development, the results are shown as yearly values rather than by seasonal variability, as done in previous sections. This approach is due to the lower frequency of systems undergoing secondary development compared to those that do not (Figure \ref{fig:filtered_lifecycle_configurations}). Additionally, as shown in Section \ref{sec:climatological_aspects}, the number of systems presenting early intensification or decay phases is comparable to those with secondary development stages. Therefore, the analysis here focuses on systems that exhibited a complete secondary development cycle: incipient, intensification, mature, decay, intensification 2, mature 2, and decay 2. Figure \ref{fig:secondary_development_combined} illustrates the track density differences when other life cycle configurations are included in the analysis. The occurrence of early intensification and decay phases impacts the overall track density distribution, with maximum density centers between \(30^\circ\)S and \(40^\circ\)S, centered near \(45^\circ\)W (Figures \ref{fig:secondary_development_combined}a and \ref{fig:secondary_development_combined}c). These centers disappear when the analysis focus only on the systems with a complete secondary development (Figures \ref{fig:secondary_development_combined}b and \ref{fig:secondary_development_combined}d).

\begin{figure}[!htbp]
\centering
\includegraphics[width=\textwidth]{figs_4/density_map_secondary_development_combined.png}
\caption[Track Density - Secondary Development Stages]{Cyclone track density for the secondary development stages. Panels A and B display the phase intensification 2, while panels C and D, phase decay 2. Panels A and C show the results when all life cycle configurations are included in the analysis, while panels B and D show the results when only the systems that exhibited a complete secondary development cycle - incipient, intensification, mature, decay, intensification 2, mature 2, and decay 2 - are analysed. The unit is cyclones per \(10^6 \, km^2\) per month.}
\label{fig:secondary_development_combined}
\end{figure}

For the second intensification phase, there is a maximum track density region south of South Africa, with \(30\) cyclones per \(10^6 \, km^2\) per month (Figure \ref{fig:intensification_2}a). This maximum density region is related to the aggregate occurrence of systems from all genesis regions and closely matches the SE-SAO cyclogenesis region reported by \citet{gramcianinov2019properties}. \citet{gramcianinov2019properties} highlighted a potential link between this region and secondary cyclogenesis. It is possible that some of the detected re-intensifications might stem from secondary cyclogenesis events or spurious links by the TRACK algorithm, which could mix weak decaying systems with incipient secondary cyclones. While secondary cyclogenesis generally follows a frontal wave post an extratropical cyclone's passage \citep{Mailier_2006, Shapiro_1997, Ford_1990, Rivals_1998}, the CycloPhaser program classifies a phase as "intensification 2" whenever the \(\zeta_{850}\) series indicates a pattern of increases in its magnitude, followed by subsequent mature and decay phases. This secondary intensification could result from a decaying wave at 850 hPa finding favorable conditions for subsequent intensification. Alternatively, the surface cyclone might have already dissipated, but there is still a signal of it in the \(\zeta_{850}\) field, causing the TRACK algorithm to continue tracking a weak decaying trough at 850 hPa. In these instances, the CycloPhaser may classify this weak intensification as 'residual' should it not progress to a mature stage (Section \ref{sec:cyclophaser_description}). Conversely, if the TRACK algorithm mistakenly associates this decaying wave with a newly forming cyclone nearby, the CycloPhaser would label this transition as "intensification 2," introducing potential analysis errors. However, distinguishing between true secondary genesis and re-intensification remains beyond this study's scope. Furthermore, notably, there is a track density maxima near Uruguay and South Brazil associated with LA-PLATA systems.


\begin{figure}[!htbp]
\centering
\includegraphics[width=\textwidth]{figs_4/density_map_intensification 2_regions_secondary.png}
\caption[Track Density - Intensification 2]{Cyclone track density during the second intensification phase for all systems with genesis in the ARG, LA-PLATA and SE-BR regions that display a complete secondary development cycle. Panel A shows the aggregate density across all regions. Panels B to D depict individual genesis regions: ARG (B), LA-PLATA (C), and SE-BR (D). The unit is cyclones per \(10^6 \, km^2\) per month.}
\label{fig:intensification_2}
\end{figure}

The second mature phase, characterized by its overall low track density, presents a diffuse spatial distribution. This low track density is primarily due to the infrequent occurrence of secondary development in SESA systems (Figure \ref{fig:filtered_lifecycle_configurations}) and the fact that it exhibits the lowest displacement and duration across all phases (Figures \ref{fig:pdf_total_distance} and \ref{fig:pdf_total_time}). However, a notable maximum track density is centered at \(60^\circ\)S and \(50^\circ\)E, with \(2\) cyclones per \(10^6 \, km^2\) per month (Figure \ref{fig:mature_2}a). This peak is mainly associated with systems from ARG and SE-BR (Figures \ref{fig:mature_2}b and \ref{fig:mature_2}d). A weaker density maximum is found to the west, between \(10^\circ\)E and \(20^\circ\)E, predominantly linked to systems from LA-PLATA (Figure \ref{fig:mature_2}c). For all regions, the track density regions are mostly displaced southeastward compared to the intensification 2 phase.


\begin{figure}[!htbp]
\centering
\includegraphics[width=\textwidth]{figs_4/density_map_mature 2_regions_secondary.png}
\caption[Track Density - Mature 2]{Cyclone track density during the second mature phase for all systems with genesis in the ARG, LA-PLATA and SE-BR regions that display a complete secondary development cycle. Panel A shows the aggregate density across all regions. Panels B to D depict individual genesis regions: ARG (B), LA-PLATA (C), and SE-BR (D). The unit is cyclones per \(10^6 \, km^2\) per month.}
\label{fig:mature_2}
\end{figure}

The second decay phase displays the systems in their southernmost and easternmost positions, with track densities extending up to \(140^\circ\)E and bordering the Antarctic continent. All three genesis regions present a track density maximum centered between \(60^\circ\)E and \(110^\circ\)E, with the resulting aggregate displaying up to \(8.5\) cyclones per \(10^6 \, km^2\) per month (Figure \ref{fig:decay_2}a). Individually, each genesis region has a track maximum of approximately \(2\) to \(3\) cyclones per \(10^6 \, km^2\) per month, with LA-PLATA displaying a second weaker maximum between \(20^\circ\)E and \(40^\circ\)E, with up to \(2\) cyclones per \(10^6 \, km^2\) per month (Figure \ref{fig:decay_2}c). This region also presents the highest density values northwestwards, contributing to the track densities near \(40^\circ\)S.


\begin{figure}[!htbp]
\centering
\includegraphics[width=\textwidth]{figs_4/density_map_decay 2_regions_secondary.png}
\caption[Track Density - Decay 2]{Cyclone track density during the second decay phase for all systems with genesis in the ARG, LA-PLATA and SE-BR regions that display a complete secondary development cycle. Panel A shows the aggregate density across all regions. Panels B to D depict individual genesis regions: ARG (B), LA-PLATA (C), and SE-BR (D). The unit is cyclones per \(10^6 \, km^2\) per month.}
\label{fig:decay_2}
\end{figure}

Across all second development phases, a weaker albeit noticeable track density region is present near the SE-BR genesis area. This region also shows occurrences of subtropical cyclones \cite{evans2012climatology, de2022future, cardoso2022synoptic}, especially during DJF months, where approximately 30\% of the systems exhibit hybrid characteristics \cite{gozzo2014subtropical}. Previous works analyzing case studies of such systems that have genesis in the SE-BR region reveal that they present complex life cycles, often transitioning between different types of cyclonic systems and presenting multiple stages of intensification, maturation, and decay \citep[e.g.,]{dias2013synoptic, veiga2008analysis, reboita2022shapiro, dutra2017structure, reboita2021iba}. Consequently, one can speculate that some of the systems related to these track density regions near SE-BR might be subtropical. However, further analysis would be needed for confirm this assertion, which is beyond the scope of the current work.


\subsection{Residual Phase}\label{residual_phase}


The residual phase, though not representing a physical stage of cyclone development, plays an important role in identifying tracking and detection inconsistencies, highlighting potential areas for methodological improvement. This phase typically includes systems that experience temporary intensification due to favorable environmental factors but fail to progress into mature cyclones. Identifying the residual phase involves detecting situations where a system's central vorticity increases outside the standard development cycle, often influenced by nearby systems or localized atmospheric conditions that promote brief intensification (Section \ref{sec:cyclophaser_description}).

Figure \ref{fig:residual} shows the track densities across all regions and seasons. A noticeable track density maximum is present near Uruguay and South Brazil, over the Brazil-Malvinas confluence \citep{gordon1989brazil}. The confluence region presents low-level baroclinicity, which can be transferred upwards in the atmosphere \citep[e.g.,]{sanders1980synoptic, vera2002cold}. Systems passing over the confluence might experience brief intensification due to favorable conditions but not achieve maturity. This could also be important for systems with genesis in SE-BR, which are often weaker \citep{sinclair1995climatology, hoskins2005new}. Furthermore, the track density maximum near \(40^\circ\)E and \(55^\circ\)S indicates that upper-level jet support \citep[e.g.,]{swart2015comparing} and a baroclinic environment \citep[e.g.,]{hoskins2005new} could similarly drive marginal intensification, leading to a residual classification.

\begin{figure}[!htbp]
\centering
\includegraphics[width=\textwidth]{figs_4/density_map_residual_aggregate.png}
\caption[Track Density - Residual Phase]{Cyclone track density during the residual phase for all systems with genesis in the ARG, LA-PLATA and SE-BR regions. The unit is cyclones per \(10^6 \, km^2\) per month.}
\label{fig:residual}
\end{figure}


\section{Summary and Concluding Remarks}

This chapter introduced CycloPhaser, a novel, open-source Python package for automated cyclone life cycle detection, available on the PyPI repository. Its low memory footprint, flexibility, and rapid deployment capabilities make it a valuable asset for climatological studies or even operational forecasts. The TRACK program identified a total of 7,931 systems from 1979 to 2020 for the cyclogenesis regions near South America. These results were narrowed down to life cycle configurations that accounted for at least 1\% of the overall life cycle types count and presented at least one mature phase. These configurations represented approximately 95\% of the total systems, amounting to 7,531 systems in total. The most common cyclone type followed a four-phase configuration: incipient, intensification, mature, and decay, accounting for about 60\% of the analyzed systems. The results from the CycloPhaser program align well with previous studies regarding cyclone behavior in the South Atlantic region, confirming known cyclogenesis regions and providing statistical insights into cyclone lifetimes and paths. The CycloPhaser program calculated an average cyclone lifetime of approximately 97.6 and 91.1 hours (4 and 3.7 days) for austral summer (DJF) and winter (JJA), respectively, and an average displacement of 4,658 and 4,777 km for DJF and JJA. The mean displacement speed ranged from 14 to 15 \(m s^{-1}\) from DJF to JJA, while the mean vorticity ranged from 4 to 4.5 \(-10^{-5} s^{-1}\). These findings largely corroborate those of previous studies for the regions and indicate a broad range of cyclone behaviors \citep{simmonds1999southern, hoskins2005new, mendes2010climatology, reboita2010regimes, gramcianinov2019properties, gramcianinov2020analysis, crespo2023assessment}.

The analysis of spatial distributions revealed distinct patterns across various phases of the cyclone life cycle:
\begin{itemize}
    \item \textbf{Incipient Phase}: The overall pattern matches the genesis regions described in the literature.
    \item \textbf{Intensification Phase}: The seasonal differences in track densities are mostly related to seasonal differences in genesis. Two maxima are present in both seasons: one mostly associated with ARG and the other with either LA-PLATA or SE-BR. However, other patterns include ARG systems being more eastward-oriented during JJA than DJF and reaching further east during DJF, while LA-PLATA systems tend to reach further southeast during JJA. Overall, the tracks are more displaced northward during JJA than DJF.
    \item \textbf{Mature Phase}: The maximum track densities are southeast of the intensification centers during DJF and east/southeast during JJA. During JJA, cyclones from the ARG region mature close to the intensification. There is a higher positioning of cyclones close to Uruguay and Southern/Southeastern Brazil during DJF than JJA.
    \item \textbf{Decay Phase}: The track density regions elongate southeastward compared to the mature phase, with a maximum density east of the Weddell Sea, mainly related to systems from ARG. Systems with an early decay phase appear as density centers in LA-PLATA. During DJF, there are more tracks positioned near Uruguay and South/Southeast Brazil than JJA.
    \item \textbf{Second Intensification Phase}: There is a high concentration of tracks south of South Africa, which might be linked to secondary cyclogenesis processes \citep{Mailier_2006, Shapiro_1997, Ford_1990, Rivals_1998}.
    \item \textbf{Second Mature Phase}: This phase generally occurs southeastward of the second intensification phase.
    \item \textbf{Second Decay Phase}: This phase occurs bordering Antarctica, predominantly east of \(40^\circ\)E.
\end{itemize}

Figure \ref{fig:density_map_Aggregate} provides an illustrative representation of the life cycle of systems with genesis in the South American cyclogenesis regions. This figure highlights the systems' path through their lifecycle, showing that after genesis, the systems follow the usual storm tracks in the South Atlantic or, in the case of systems with genesis in ARG and LA-PLATA, the subtropical storm path branch. Although systems from ARG and SE-BR tend to have their genesis and intensification over the ocean, LA-PLATA systems tend to start intensifying close to the continent, which can impact coastal communities \citep{de2021ocean}. Moreover, systems from SE-BR and LA-PLATA have their mature phases — where they are most intense and, therefore, transfer more momentum from the atmosphere to the ocean — near oceanic regions close to South and Southeast Brazil. These regions host large urban centers and significant economic activities, such as oil drilling.

\begin{figure}[!htbp]
\centering
\includegraphics[width=\textwidth]{figs_4/density_map_Aggregate.png}
\caption[Illustrative Representation of Cyclone Life Cycle]{Track densities normalized for each region, using yearly values for all phases. For computing the track densities for the secondary development phases, only systems that display a complete secondary development cycle were used. Contours representing normalized track densities above 0.8 were plotted for each region individually.}
\label{fig:density_map_Aggregate}
\end{figure}

The results presented here contribute significantly to the climatology of South Atlantic cyclones. CycloPhaser's detailed analysis of cyclone life cycle types, particularly focusing on configurations accounting for significant proportions of cyclone activity, enhances our understanding of the spatial distribution of cyclone phases. This analysis reaffirms the common four-phase configuration (incipient, intensification, mature, and decay) and sheds light on variations in cyclone development cycles, contributing to a better understanding of how different regions within the South Atlantic influence the formation and evolution of distinct cyclone types. These insights are valuable for improving weather forecasting and developing effective mitigation strategies in the region.

Despite the valuable contributions of the CycloPhaser program, it is important to acknowledge the study's limitations. The analysis at a single atmospheric level offers a concentrated but partial view of the cyclone life cycle. Consequently, our conclusions about cyclone characteristics throughout their life cycle predominantly reflect the processes observable at the 850 hPa level. This methodological choice inherently limits our ability to fully explore the cyclones' vertical development and the complex interactions occurring at different atmospheric altitudes. Future research could benefit from incorporating a multi-level approach, providing a more holistic understanding of cyclone dynamics across the troposphere.

One of the unique contributions of this research is the assessment of cyclone life cycle stages along their storm tracks, a perspective not commonly explored in existing cyclone track climatologies. This first application has highlighted the complexity of cyclone systems and opened several avenues for future research. The approach presented here allows for the analysis of distinct dynamical mechanisms linked to the development of cyclonic systems at different stages of their life cycle. Investigating cyclonic life cycles in climate projections might be particularly relevant for coastal management in South America, posing questions such as: Will future cyclones intensify and mature closer to or further from coastal regions? Will these phases last longer? Furthermore, using the method presented here, future studies can delve deeper into the primary processes associated with each development stage in cyclone life cycles, including the secondary development stages reported herein, and determine whether some of these cases represent secondary genesis.



